%! AP Statistics — Unit 2, Lesson 2.6 — Homework — Answer Key
\documentclass[10pt]{article}
\usepackage{apstats-article}
\usepackage{apstats-key}

\begin{document}

\pageheader{Unit 2, Lesson 2.6}{Homework --- Answer Key}
\namedateperiod

\vspace{0.18in}

% ── DIRECTIONS ───────────────────────────────────────────────────────────────
\noindent\textbf{Directions:} Show all work. Write all interpretations in complete sentences
with context. Flag any extrapolation.

\vspace{0.14in}

% ── PROBLEM 1 ─────────────────────────────────────────────────────────────────
\noindent\textbf{1.}\enspace A used-car dealer tracked mileage (in thousands of miles, $x$)
and asking price (in \$1000s, $y$) for 18 vehicles. The data ranged from \textbf{8 to 95
thousand miles}. Technology produced:
\[
  \widehat{\text{price}} = 32.0 - 0.24(\text{mileage})
\]

\begin{enumerate}[label=\textbf{(\alph*)}, leftmargin=2em, itemsep=0.18in]

  \item Interpret the \textbf{slope} in context.\\[8pt]
        \ansline{For each additional thousand miles of mileage, the predicted asking price
        decreases by approximately \$240.}

  \item Interpret the \textbf{$y$-intercept} in context. Is it meaningful? Explain.\\[8pt]
        \ansline{When mileage = 0 (a brand-new car with no miles), the predicted asking price
        is \$32{,}000. This is \textbf{not meaningful in context} because 0 miles is outside
        the data range (8--95 thousand miles) --- the model was built from used cars only.}

  \item Predict the asking price of a car with \textbf{50 thousand miles}. Show your work and
        classify as interpolation or extrapolation.\\[8pt]
        \ansline{$\widehat{\text{price}} = 32.0 - 0.24(50) = 32.0 - 12.0 = \$20{,}000$.\quad
        \textbf{Interpolation} (50 is within 8--95), so this prediction is reliable.}

  \item Predict the asking price of a car with \textbf{120 thousand miles}. Show your work.
        Is this prediction reliable? Why or why not?\\[8pt]
        \ansline{$\widehat{\text{price}} = 32.0 - 0.24(120) = 32.0 - 28.8 = \$3{,}200$.
        This is \textbf{extrapolation} (120 > 95), so the prediction is unreliable. The linear
        pattern may not continue beyond the data range.}

\end{enumerate}

\vspace{0.18in}

% ── PROBLEM 2 ─────────────────────────────────────────────────────────────────
\noindent\textbf{2.}\enspace A sports scientist measured arm span (cm, $x$) and height (cm,
$y$) for 40 swimmers. Arm spans ranged from \textbf{160 cm to 210 cm}. The regression
equation is:
\[
  \widehat{\text{height}} = 12.0 + 0.95(\text{arm span})
\]

\begin{enumerate}[label=\textbf{(\alph*)}, leftmargin=2em, itemsep=0.18in]

  \item Write a complete interpretation of the \textbf{slope} in context.\\[8pt]
        \ansline{For each additional centimeter of arm span, the predicted height increases by
        approximately 0.95 cm.}

  \item Interpret the \textbf{$y$-intercept}. Is it meaningful in this context?\\[8pt]
        \ansline{When arm span = 0 cm, the predicted height is 12.0 cm. This is \textbf{not
        meaningful}: no swimmer has a 0 cm arm span, 0 is far outside the data range
        (160--210 cm), and a height of 12 cm is nonsensical for an adult swimmer. The
        intercept is a mathematical artifact of fitting the line.}

  \item Predict the height of a swimmer with an arm span of \textbf{185 cm}. Show work.\\[8pt]
        \ansline{$\widehat{\text{height}} = 12.0 + 0.95(185) = 12.0 + 175.75 = 187.75$ cm.
        This is interpolation (185 cm is within 160--210 cm) and is reliable.}

\end{enumerate}

\vspace{0.18in}

% ── PROBLEM 3 ─────────────────────────────────────────────────────────────────
\noindent\textbf{3.}\enspace Computer output from a regression of $y$ = plant height (cm) on
$x$ = weeks of growth is shown below.

\vspace{8pt}
\renewcommand{\arraystretch}{1.4}
\begin{tabularx}{\linewidth}{@{} l X X X X @{}}
\rowcolor{sky}
\textbf{Term} & \textbf{Coef} & \textbf{SE Coef} & \textbf{T-Value} & \textbf{P-Value} \\
\midrule
Constant       & 2.4  & 0.8  & 3.0  & 0.008 \\
Weeks of growth & 3.7  & 0.3  & 12.3 & 0.000 \\
\end{tabularx}

\vspace{10pt}

\begin{enumerate}[label=\textbf{(\alph*)}, leftmargin=2em, itemsep=0.18in]

  \item Write the regression equation using variable names.\\[8pt]
        \ansline{$\widehat{\text{height}} = 2.4 + 3.7(\text{weeks of growth})$}

  \item Interpret the \textbf{slope} in context.\\[8pt]
        \ansline{For each additional week of growth, the predicted plant height increases
        by approximately 3.7 cm.}

  \item Predict the height of a plant after \textbf{5 weeks} of growth.
        Show work.\\[8pt]
        \ansline{$\widehat{\text{height}} = 2.4 + 3.7(5) = 2.4 + 18.5 = 20.9$ cm.}

\end{enumerate}

\vspace{0.18in}

% ── EXTENSION ─────────────────────────────────────────────────────────────────
\begin{extensionbox}
\small\textbf{Extension --- Optional}

Find a published regression equation in a science article, textbook, or online database.
Write out: (1) the equation with variable names, (2) an interpretation of the slope in
context, (3) whether the $y$-intercept is meaningful, and (4) one prediction the equation
could make and whether it is interpolation or extrapolation. Attach your source.
\end{extensionbox}

\vspace{0.16in}

% ── PREVIEW ───────────────────────────────────────────────────────────────────
\begin{tcolorbox}[colback=greenbg, colframe=greenacc, boxrule=0.8pt, arc=2mm,
    left=4mm, right=4mm, top=3mm, bottom=3mm]
\small\textbf{Preview --- Lesson 2.7}\\[4pt]
The regression equation gives a \emph{predicted} value $\hat{y}$ for every $x$. But the
actual $y$ value for each data point is rarely exactly $\hat{y}$. In Lesson 2.7 we measure
that discrepancy:
\[
  \text{residual} = y - \hat{y}
\]
\textbf{Think ahead:} For the car example, if a 5-year-old car actually sold for \$20{,}000,
what is its residual? Is the prediction an overestimate or an underestimate?\\[6pt]
\ansline{$\hat{y} = 19.15$ thousand; residual $= 20.0 - 19.15 = 0.85$ thousand (\$850).
Positive residual $\Rightarrow$ prediction was an \textbf{underestimate}.}
\end{tcolorbox}

\begin{teachernote}
Problem 1(b): students may say the intercept \emph{is} meaningful because ``a new car has
0 miles.'' Accept if they clearly state 0 is outside the data range and the model was built
from used cars --- the key AP skill is acknowledging the range limitation, not just whether
0 is physically possible. Problem 3 introduces reading a computer output table; preview
Lesson 2.9 if students ask about SE Coef and P-Value columns.
\end{teachernote}

\end{document}
